\chapter{The beer domain}

The following sections present the analysis of the domain over which BeerEX
reasons, the dimensions along which beer is described, the taxonomy adopted to
organize the styles, and the principles that govern the association between beer
and food. A brief introduction to the notion of a beer \emph{style} is followed
by the distinction between styles and commercial brands; the measurable and
perceptual parameters of a style are then defined and tabulated; the styles are
classified into families; the three mechanisms of food pairing are described; and,
finally, the documentary sources and the reproducible extraction pipeline that
keep the knowledge base faithful to those sources are presented.

\section{Beer styles versus brands}

The domain modelled by BeerEX is that of beer \emph{styles} rather than that of
individual commercial \emph{brands}. A style is an abstract, bounded category
---``German-Style Pilsener'', ``Irish-Style Dry Stout'', ``Belgian-Style
Tripel''--- defined by a characteristic range of measurable parameters and a
recognizable sensory profile, and codified by an authoritative body. A brand, by
contrast, is a single product of a single brewery: an instance of a style, often
itself drifting from batch to batch, and numbering in the tens of thousands
worldwide. Choosing the style as the unit of the domain has three advantages.
First, it is \emph{finite and stable}: the catalogue of recognized styles changes
slowly and is published, whereas the catalogue of brands is unbounded and
volatile. Second, it is \emph{described by data}: each style is associated with
numeric ranges (alcohol, bitterness, colour, gravity, carbonation) and with a
qualitative profile, so that the rules of the system have a firm, quantifiable
substrate on which to operate. Third, it is the \emph{level at which food pairing
is meaningful}: the pairing guides associate dishes with styles, not with
particular labels, because it is the style's flavour architecture, and not a
brand's marketing, that determines whether a beer complements or clashes with a
dish.

The original version of BeerEX encoded only a handful of symbolic attributes per
style (alcohol, colour, flavour, fermentation, carbonation). The current
knowledge base, defined in \code{clips/beer-styles.fct}, retains that symbolic
backbone but enumerates \textbf{79} styles, each assigned to one of fifteen
families and described by the symbolic slots \code{alcohol}, \code{color},
\code{flavor}, \code{fermentation}, \code{carbonation} and a reference
\code{link} to the CraftBeer.com style page. This symbolic dataset is then
enriched, style by style, by two further fact bases extracted directly from the
source documents: the numeric ranges of \code{clips/beer-ranges.clp}
(\textbf{73} of the 79 styles carry characteristic ranges; the remaining six are
catch-all specialty or wild categories whose values span the entire scale and are
therefore not informative) and the comprehensive qualitative, serving and pairing
profiles of \code{clips/beer-data.clp} (\textbf{79} profiles and \textbf{73}
food pairings). Throughout this report, ``style'' and ``beer'' are used
interchangeably to denote one such entry.

\section{Beer parameters}

Each style is described along two complementary kinds of dimension: a set of
\emph{quantitative} parameters, which are continuous physical or chemical
quantities measured during brewing or analysis, and a set of \emph{qualitative}
parameters, which are the perceptual descriptors recorded by tasters. To these is
added the \emph{serving} information ---temperature and glassware--- that
completes the practical profile of a style.

\subsection{Quantitative parameters}

The quantitative parameters are the seven continuous quantities listed in
Table~\ref{tab:quant}. They are extracted as per-style closed intervals
(a lower and an upper bound) from the source guide and stored in the
\code{beer-range} deftemplate as the multislots \code{abv}, \code{ibu},
\code{srm}, \code{og}, \code{fg} and \code{co2}. The three of them that also
serve as the system's principal linguistic axes ---\emph{colour} (SRM),
\emph{alcohol} (ABV) and \emph{bitterness} (IBU)--- are the ones along which the
user expresses a preference and along which the fuzzy backend builds its
membership functions.

\begin{itemize}[leftmargin=2em]
  \item \textbf{Alcohol by volume (ABV)} is the fraction of the beverage's volume
        that is ethanol, the most familiar measure of strength.
  \item \textbf{International Bitterness Units (IBU)} quantify bitterness by the
        concentration of isomerized alpha-acids derived from the hops.
  \item \textbf{Standard Reference Method (SRM)} measures colour by light
        absorbance: low values denote pale, straw-coloured beers and high values
        deep brown to black ones.
  \item \textbf{Original gravity (OG)} is the density of the wort relative to
        water before fermentation; it reflects the amount of fermentable sugar
        available and hence the potential strength and body.
  \item \textbf{Final gravity (FG)} is the density after fermentation; the
        residual sugar it represents contributes to sweetness and body.
  \item \textbf{Apparent attenuation} is the proportion of the original extract
        consumed by the yeast, derived from OG and FG as
        $(\mathrm{OG}-\mathrm{FG})/(\mathrm{OG}-1)$; it expresses how dry or how
        sweet the finished beer is.
  \item \textbf{Carbonation (CO\textsubscript{2} volumes)} is the amount of
        dissolved carbon dioxide, measured in volumes of gas per volume of beer;
        it governs the perceived effervescence.
\end{itemize}

Table~\ref{tab:quant} also reports, for each parameter, the overall extent
actually spanned by the dataset, that is, the lowest lower bound and the highest
upper bound observed across the styles that carry that parameter in
\code{beer-ranges.clp}.

\begin{longtable}{@{}p{0.26\textwidth}p{0.40\textwidth}p{0.12\textwidth}p{0.14\textwidth}@{}}
  \caption{The quantitative parameters of a beer style, with the overall range
  spanned by the dataset.}\label{tab:quant}\\
  \toprule
  \textbf{Parameter} & \textbf{What it measures} & \textbf{Unit} &
  \textbf{Dataset range} \\
  \midrule
  \endfirsthead
  \multicolumn{4}{@{}l}{\small\itshape Table~\ref{tab:quant} (continued)}\\
  \toprule
  \textbf{Parameter} & \textbf{What it measures} & \textbf{Unit} &
  \textbf{Dataset range} \\
  \midrule
  \endhead
  \midrule
  \multicolumn{4}{r}{\small\itshape continued on next page}\\
  \endfoot
  \bottomrule
  \endlastfoot
  Alcohol by volume (ABV) &
    ethanol content, i.e.\ strength &
    \% vol & 2.8--14.2 \\
  Bitterness (IBU) &
    isomerized alpha-acid concentration from hops &
    IBU & 3--100 \\
  Colour (SRM) &
    colour by light absorbance, pale to black &
    SRM & 2--50 \\
  Original gravity (OG) &
    wort density before fermentation (fermentable sugar) &
    s.g. & 1.028--1.140 \\
  Final gravity (FG) &
    density after fermentation (residual sugar) &
    s.g. & 1.000--1.032 \\
  Apparent attenuation &
    fraction of extract fermented, derived from OG and FG &
    \% & derived \\
  Carbonation &
    dissolved CO\textsubscript{2}, perceived effervescence &
    CO\textsubscript{2} vol. & 1.0--4.5 \\
\end{longtable}

\subsection{Qualitative parameters}

The qualitative parameters are the perceptual descriptors recorded, per style, in
the \code{beer-profile} deftemplate of \code{clips/beer-data.clp}. Unlike the
numeric ranges, they are free-text sensory characterizations taken verbatim from
the source guide. The slots are the following.

\begin{itemize}[leftmargin=2em]
  \item \textbf{Perceived alcohol} (\code{alcohol}): how evident the alcohol is on
        the palate, on a scale running through ``Not Detectable'', ``Mild'',
        ``Noticeable'', ``Hot'' and ``Harsh'' (e.g.\ ``Not Detectable to Mild''
        for a Bohemian-Style Pilsener, ``Noticeable to Hot'' for an American
        Barley Wine).
  \item \textbf{Colour} (\code{colour}): the visual colour described in words
        (e.g.\ ``Straw to Pale'', ``Copper to Reddish Brown'', ``Black''),
        complementing the numeric SRM range.
  \item \textbf{Clarity} (\code{clarity}): the transparency of the beer, from
        ``Brilliant'' and ``Clear'' through ``Slight Haze'' to ``Hazy'' and
        ``Opaque''.
  \item \textbf{Body} (\code{body}): the palate weight or mouthfeel, from
        ``Drying'' and ``Soft'' to ``Moderate'', ``Mouth-Coating'' and
        ``Sticky''.
  \item \textbf{Palate carbonation} (\code{palate-carbonation}): the perceived
        effervescence on the palate (e.g.\ ``Low'', ``Medium'', ``Medium to
        High''), the sensory counterpart of the numeric CO\textsubscript{2}
        parameter.
  \item \textbf{Finish} (\code{finish}): the length and character of the
        aftertaste, from ``Short'' to ``Long''.
  \item \textbf{Hop aroma and flavour} (\code{hop}): the contribution of the hops
        to aroma, flavour and bitterness.
  \item \textbf{Malt aroma and flavour} (\code{malt}): the contribution of the
        malt, including the characteristic descriptors of the style (caramel,
        toast, biscuit, chocolate, coffee, and so on).
\end{itemize}

The yeast character, finally, is captured indirectly through the symbolic
\code{fermentation} slot of \code{beer-styles.fct}, which records whether the
style is \emph{top}-fermented (ale), \emph{bottom}-fermented (lager),
spontaneously or mixed (\emph{wild}), or a combination thereof. The profile also
records the style's \code{country} of origin.

\subsection{Serving}

Two further slots of \code{beer-profile} describe how a style is meant to be
served, and so complete its practical profile.

\begin{itemize}[leftmargin=2em]
  \item \textbf{Temperature} (\code{temperature}): the recommended serving
        temperature, given as a Fahrenheit range. Lighter styles are served
        cooler ---several pale and wheat styles at ``40--45\,\textdegree
        F''--- and stronger, darker styles warmer (typically ``50--55\,\textdegree
        F''), so that their aromatics open up.
  \item \textbf{Glassware} (\code{glass}): the glass that best presents the style,
        drawn from a small vocabulary of vessels ---among them the Tulip, the
        Nonic Pint, the Flute, the Snifter, the Goblet, the Vase and the
        Thistle--- each suited to the style's strength, carbonation and aromatic
        intensity.
\end{itemize}

\section{Beer style families}

The 79 styles are organized into the \textbf{fifteen} families recorded in the
\code{style} slot of \code{beer-styles.fct}. The classification follows the
CraftBeer.com grouping, which combines the fundamental fermentation distinction
(ale versus lager) with colour and strength: the families range from the pale,
lightly fermented \emph{Pilsener \& Pale Lager} and \emph{Pale Ale} groups,
through the malt-forward \emph{Brown Ale}, \emph{Dark Lager}, \emph{Bock},
\emph{Porter} and \emph{Stout} groups, to the high-strength \emph{Strong Ale} and
the aromatic \emph{Belgian Style} and \emph{Wild/Sour} groups, with the
\emph{India Pale Ale}, \emph{Wheat Beer}, \emph{Scottish-Style Ale} and
\emph{Hybrid Beer} families and a residual \emph{Specialty Beer} family that
collects flavoured, barrel-aged and otherwise non-canonical beers.
Table~\ref{tab:families} lists the fifteen families, their cardinality in the
dataset, and a representative style for each, with that style's symbolic
descriptors and the numeric ABV, IBU and SRM ranges actually recorded for it.

\begin{longtable}{@{}p{0.18\textwidth}p{0.04\textwidth}p{0.24\textwidth}p{0.14\textwidth}p{0.09\textwidth}p{0.09\textwidth}p{0.07\textwidth}@{}}
  \caption{The fifteen beer families, with a representative style and its recorded
  ABV (\% vol), IBU and SRM ranges.}\label{tab:families}\\
  \toprule
  \textbf{Family} & \textbf{\#} & \textbf{Representative style} &
  \textbf{Colour / flavour} & \textbf{ABV} & \textbf{IBU} & \textbf{SRM} \\
  \midrule
  \endfirsthead
  \multicolumn{7}{@{}l}{\small\itshape Table~\ref{tab:families} (continued)}\\
  \toprule
  \textbf{Family} & \textbf{\#} & \textbf{Representative style} &
  \textbf{Colour / flavour} & \textbf{ABV} & \textbf{IBU} & \textbf{SRM} \\
  \midrule
  \endhead
  \midrule
  \multicolumn{7}{r}{\small\itshape continued on next page}\\
  \endfoot
  \bottomrule
  \endlastfoot
  Pale Ale & 5 & American Pale Ale & amber, hoppy-bitter & 4.4--5.4 & 30--50 & 6--14 \\
  India Pale Ale & 3 & American IPA & amber, hoppy-bitter & 6.3--7.6 & 50--70 & 6--14 \\
  Pilsener \& Pale Lager & 5 & German-Style Pilsener & pale, crisp-clean & 4.6--5.3 & 25--40 & 3--4 \\
  Dark Lager & 5 & German-Style Dunkel & brown, dark-roasty & 4.8--5.3 & 16--25 & 15--17 \\
  Brown Ale & 3 & American Brown Ale & brown, malty-sweet & 4.2--6.3 & 25--45 & 15--26 \\
  Wheat Beer & 6 & German-Style Hefeweizen & pale, fruity-spicy & 4.9--5.6 & 10--15 & 3--9 \\
  Belgian Style & 7 & Belgian-Style Tripel & pale/amber, fruity-spicy & 7.1--10.1 & 20--45 & 4--9 \\
  Strong Ale & 4 & American Barley Wine & brown, hoppy-bitter & 8.5--12.2 & 60--100 & 11--18 \\
  Bock & 4 & German-Style Doppelbock & brown/dark, malty-sweet & 6.6--7.9 & 17--27 & 12--30 \\
  Porter & 5 & Robust Porter & dark, dark-roasty & 5.1--6.6 & 25--40 & 30--50 \\
  Stout & 5 & Irish-Style Dry Stout & dark, dark-roasty & 4.2--5.3 & 30--40 & 40--50 \\
  Scottish-Style Ale & 2 & Scotch Ale/Wee Heavy & brown/dark, malty-sweet & 6.6--8.5 & 25--35 & 15--30 \\
  Hybrid Beer & 6 & California Common & amber/brown, hoppy-bitter & 4.6--5.7 & 35--45 & 8--15 \\
  Wild/Sour & 6 & Contemporary Gose & pale, sour-tart-funky & 4.4--5.4 & 10--15 & 3--3 \\
  Specialty Beer & 13 & American Black Ale & dark, hoppy-bitter & 6.3--7.6 & 50--70 & 35--50 \\
\end{longtable}

The table makes the structuring role of the three linguistic axes apparent.
Colour (SRM) separates the pale families (Pilsener, Pale Ale, Wheat) from the
dark ones (Porter, Stout) almost completely; alcohol (ABV) singles out the Strong
Ale and Belgian families; and bitterness (IBU) distinguishes the hop-forward
India Pale Ale and Strong Ale families from the malt-forward Bock, Dark Lager and
Brown Ale families. These are precisely the gradients along which BeerEX's fuzzy
membership functions are calibrated.

\section{Food pairing}

Beer-and-food pairing, as codified by the Brewers Association, is not arbitrary
but rests on three complementary mechanisms, which the system encodes as the
principles by which a dish is matched to a style.

\begin{itemize}[leftmargin=2em]
  \item \textbf{Complement.} Shared or harmonious flavours reinforce one another,
        as when the roasted, coffee-like notes of a stout echo those of a
        chocolate dessert.
  \item \textbf{Contrast.} Opposing flavours balance one another, as when the
        sweetness of a malty beer offsets a savoury or spicy dish.
  \item \textbf{Cut.} The carbonation and bitterness of the beer cleanse the
        palate of fat and intensity, refreshing it between bites of a rich dish.
\end{itemize}

On the basis of these mechanisms, the source guides associate each style with the
\emph{cheese}, \emph{entrée} and \emph{dessert} that best match it. BeerEX records
these associations in the \code{pairing} deftemplate of
\code{clips/beer-data.clp}, with \textbf{73} of the styles carrying a triple of
authoritative pairings. The system can therefore reason from a meal or dish back
to a style: a request centred on a particular cheese, a particular entrée
(roasted or grilled meats, seafood, sausages, spicy dishes) or a dessert is
matched against these triples, and the matching styles are scored together with
the measured attributes. Representative pairings drawn directly from the
knowledge base are shown in Table~\ref{tab:pairings}.

\begin{longtable}{@{}p{0.27\textwidth}p{0.20\textwidth}p{0.24\textwidth}p{0.21\textwidth}@{}}
  \caption{Representative food pairings, one per family, from
  \code{beer-data.clp}.}\label{tab:pairings}\\
  \toprule
  \textbf{Style} & \textbf{Cheese} & \textbf{Entrée} & \textbf{Dessert} \\
  \midrule
  \endfirsthead
  \multicolumn{4}{@{}l}{\small\itshape Table~\ref{tab:pairings} (continued)}\\
  \toprule
  \textbf{Style} & \textbf{Cheese} & \textbf{Entrée} & \textbf{Dessert} \\
  \midrule
  \endhead
  \midrule
  \multicolumn{4}{r}{\small\itshape continued on next page}\\
  \endfoot
  \bottomrule
  \endlastfoot
  American Pale Ale & Mild or Medium Cheddar & Roasted or Grilled Meats & Apple Pie \\
  American IPA & Blue Cheeses & Spicy Tuna Roll & Persimmon Rice Pudding \\
  German-Style Pilsener & White Cheddar & Salads & Shortbread Cookies \\
  German-Style Dunkel & Washed-Rind Munster & Sausages & Candied Ginger Beer Cake \\
  American Brown Ale & Aged Gouda & Vegetables & Pear Fritters \\
  German-Style Hefeweizen & Ch\`evre & Seafood & Key Lime Pie \\
  Belgian-Style Tripel & Triple Creme & Roasted Turkey & Caramelized Banana Creme Brulee \\
  American Barley Wine & Strong Blue Cheeses & Beef Cheek & Rich Desserts \\
  German-Style Doppelbock & Strong Cheeses & Ham & German Chocolate Cake \\
  Robust Porter & Gruy\`ere & Roasted or Grilled Meats & Chocolate Peanut Butter Cookies \\
  Irish-Style Dry Stout & Irish Cheddar & Ham & Chocolate Desserts \\
  Scotch Ale/Wee Heavy & Pungent Cheeses & Game & Creamy Desserts with Fruit \\
  California Common & Feta & Pork Loin & Bread Pudding \\
  Contemporary Gose & Queso Fresco & Watermelon Salad & Greek Yogurt Lemon Mousse \\
  American Black Ale & Blue Cheeses and Aged Gouda & Grits & Chocolate Truffles \\
\end{longtable}

The pairings illustrate the three mechanisms at work: the roasted Irish-Style Dry
Stout matched to chocolate desserts \emph{complements} like with like; the malty
German-Style Doppelbock matched to ham \emph{contrasts} sweet against savoury;
and the highly carbonated, bitter American IPA matched to a spicy tuna roll
\emph{cuts} through heat and fat. The same data also support meal- and
dish-oriented reasoning of the kind a user typically begins from ---for instance
pizza, or a cheese course--- since dishes such as roasted or grilled meats,
seafood, sausages and the various cheeses recur across families and can be used
as entry points into the style space.

\section{Sources and extraction}

The knowledge base is built entirely from authoritative material published by
CraftBeer.com and the Brewers Association, so that every fact in the system is
traceable to a published source.

\begin{itemize}[leftmargin=2em]
  \item the \emph{2017 Beer Styles Study Guide}, which provides, per style, the
        numeric ranges (ABV, IBU, SRM, OG, FG, CO\textsubscript{2} volumes) and
        the full qualitative and serving profile;
  \item the \emph{Beer \& Food Course} and the accompanying food-pairing chart,
        which provide the three pairing mechanisms (complement, contrast, cut)
        and the authoritative cheese / entrée / dessert associations.
\end{itemize}

Rather than transcribing this material by hand, BeerEX mines it with two
reproducible scripts, so that the knowledge base stays faithful to ---and is
regenerable from--- the documents.

\begin{itemize}[leftmargin=2em]
  \item \code{scripts/extract\_ranges.py} parses the Study Guide PDF (via
        \code{pdftotext}), locates each style page, reads off the six numeric
        ranges, and matches each guide entry against the 79 styles of
        \code{beer-styles.fct} by fuzzy string matching. It deliberately discards
        catch-all spans ---ranges too wide to be characteristic, as listed by the
        specialty and wild categories--- so that only informative ranges enter
        the engine. It writes the auditable table \code{data/beer-style-ranges.csv}
        and the CLIPS facts \code{clips/beer-ranges.clp}.
  \item \code{scripts/extract\_kb.py} parses both the Study Guide PDF, for the
        eleven qualitative, serving and origin fields, and the pairings
        spreadsheet, for the cheese / entrée / dessert triples, again aligning
        every record to the 79 styles. It writes the comprehensive table
        \code{data/beer-styles.csv} and the CLIPS facts \code{clips/beer-data.clp}
        (the \code{beer-profile} and \code{pairing} deftemplates and their
        deffacts).
\end{itemize}

Because the source documents are kept outside version control, the generated CSV
tables and \code{.clp} fact bases serve as the auditable, committed record of the
extraction; re-running the scripts against the same documents reproduces them
exactly. This pipeline is what allows BeerEX to replace the hand-assigned weights
of its first version with a data-grounded knowledge base, and it is on this
foundation that the certainty-factor rules and the fuzzy backend described in the
following chapters are built.
